\chapter{Why Chord--Tangent Equals Torus Addition}
\label{ch:transport}

We now put the two laws next to each other.  The analytic law is
\[
  \Phi(u)\oplus\Phi(v)=\Phi(u+v).
\]
The geometric law is the chord--tangent operation \(\boxplus\).  The main theorem
says these are the same.

\begin{theorem}[Transport theorem]
\label{thm:transport-new}
For all \(P,Q\in E(\C)\),
\[
  P\oplus Q=P\boxplus Q.
\]
\end{theorem}

\begin{proof}
Write \(P=\Phi(u)\), \(Q=\Phi(v)\), and let \(R\) be the third intersection of the
line through \(P\) and \(Q\).  By \cref{prop:line-divisor},
\([P]+[Q]+[R]-3[O]\) is principal.  Abel's theorem therefore gives
\[
  u+v+r\equiv 0 \pmod{\Lambda},
\]
where \(R=\Phi(r)\).  Hence \(R=\Phi(-(u+v))\).  Since \(\wp\) is even and
\(\wp'\) is odd, reflecting \(R\) gives \(\Phi(u+v)\).  Thus the chord--tangent
result is exactly the transported torus sum.
\end{proof}

\begin{corollary}
The chord--tangent law is associative and commutative.
\end{corollary}

\begin{proof}
It is transported from the associative and commutative law on \(\C/\Lambda\).
\end{proof}

\begin{keyidea}{No arbitrary convention}
The reflection is forced by the equation \(P+Q+R=O\).  The third intersection is
not the sum; its inverse is the sum.
\end{keyidea}

\section*{Exercises}
\addcontentsline{toc}{section}{Exercises}

\exerciseintent{Locate the exact Abel-sum calculation inside the transport proof.}
\begin{exercise}
In the proof of \cref{thm:transport-new}, write the line divisor
\([P]+[Q]+[R]-3[O]\), choose torus preimages \(u,v,r,0\), and compute its Abel
sum.  Identify the exact equality where Abel's theorem is used.
\end{exercise}

\exerciseintent{Compare direct cubic associativity with linear associativity on the torus.}
\begin{exercise}
List the third-intersection computations needed to compare
\[
  (P\boxplus Q)\boxplus R
  \quad\text{and}\quad
  P\boxplus(Q\boxplus R)
\]
directly on the cubic.  Then write the corresponding torus calculation
\[
  (u+v)+w=u+(v+w)
\]
and explain which side is linear.
\end{exercise}

\exerciseintent{Verify the identity element using the projective vertical line \(X=xZ\).}
\begin{exercise}
For \(P=(x,y)\), compute the three intersections of the projective line
\(X=xZ\) with the cubic.  Check that this line contains \(P\) and
\(O=[0:1:0]\).  Use the reflection step to show that \(O\boxplus P=P\).
\end{exercise}

\exerciseintent{Derive \(P\boxplus Q=\Phi(u+v)\) from the line divisor and Abel sum.}
\begin{exercise}
Let \(P=\Phi(u)\) and \(Q=\Phi(v)\).  Give a step-by-step proof that the
chord--tangent construction gives \(\Phi(u+v)\).  Your proof must include the
Abel-sum equation \(u+v+r\equiv0\pmod{\Lambda}\) and the computation
\(\Phi(-r)=\Phi(u+v)\).
\end{exercise}

\exerciseintent{Identify the reflection step as the torus operation \(r\mapsto -r\).}
\begin{exercise}
In the previous exercise, compute \(r=-(u+v)\).  Then use
\[
  \Phi(-r)=\ominus\Phi(r)
\]
to identify the exact point where the reflection \(\ominus R\) enters.
\end{exercise}
